\textbf{Program 1: One binary tree, two complete implementations.}

\textbf{File:} \texttt{binary\_tree.cpp}.

\textbf{Chapter connection:} Recursive traversal code; Array representation of a binary tree.

Translate left/right pointers into indexes 2i+1 and 2i+2. Compare shape and all traversals after update, insertion, and general-tree deletion. EMPTY is reserved for unused slots in this positive-valued dataset. Destroy linked nodes in postorder.

\begin{adslabfigure}[How this program runs]
\begin{adslabgraphic}
\begin{tikzpicture}[font=\small,execute at begin node={\hyphenpenalty=10000\exhyphenpenalty=10000}]
\node[utpstep,text width=2.50cm,minimum height=1.15cm] (s0) at (0.0,0) {\textbf{1}\\Nodes and array};
\node[utpstep,text width=2.50cm,minimum height=1.15cm] (s1) at (2.95,0) {\textbf{2}\\Three traversals};
\node[utpstep,text width=2.50cm,minimum height=1.15cm] (s2) at (5.9,0) {\textbf{3}\\Update / insert / delete};
\node[utpstep,text width=2.50cm,minimum height=1.15cm] (s3) at (8.85,0) {\textbf{4}\\Compare and release};
\draw[utparrow] (s0.east)--(s1.west);
\draw[utparrow] (s1.east)--(s2.west);
\draw[utparrow] (s2.east)--(s3.west);
\end{tikzpicture}
\end{adslabgraphic}
\end{adslabfigure}


\begin{utpsteps}\setcounter{enumi}{0}\item \textbf{Headers and initial definitions.}\end{utpsteps}

\texttt{binary\_tree.cpp} --- lines 14--30.

\lstinputlisting[firstline=14,lastline=30]{all_experiments/Lab07/binary_tree.cpp}


\tcbbreak

\begin{utpsteps}\setcounter{enumi}{1}\item \textbf{Implement inorder, preorder, postorder and helpers.}\end{utpsteps}

\texttt{binary\_tree.cpp} --- lines 31--62.

% Copied from chapters/07-trees-i.tex (fig:traversals).
% Keep the chapter drawing, type sizes, colours, and node styles unchanged.
\begin{utpcompactlabfigure}[Inorder, preorder and postorder compared]
\begin{labgraphic}
\begin{tikzpicture}[font=\small]
\newcommand{\ttree}[5]{%
 % #1 xshift, #2 title, #3..#5 = orders for A..G given as csv
 \begin{scope}[xshift=#1 cm]
  \node[font=\bfseries\color{UTPNavy}] at (0,1.35) {#2};
  \node[utpnode,minimum size=7.5mm] (n1) at (0,0.55) {A};
  \node[utpnode,minimum size=7.5mm] (n2) at (-0.95,-0.45) {B};
  \node[utpnode,minimum size=7.5mm] (n3) at (0.95,-0.45) {C};
  \node[utpnode,minimum size=7.5mm] (n4) at (-1.5,-1.45) {D};
  \node[utpnode,minimum size=7.5mm] (n5) at (-0.4,-1.45) {E};
  \draw[utpedge] (n1)--(n2); \draw[utpedge] (n1)--(n3);
  \draw[utpedge] (n2)--(n4); \draw[utpedge] (n2)--(n5);
  \foreach \nd/\ord in {#3}{
    \node[circle,fill=UTPOrange,text=white,font=\tiny\bfseries,inner sep=1.1pt]
      at ($(\nd.north east)+(0.055,0.055)$) {\ord};
  }
  \node[utplabel,anchor=north,align=center,color=UTPMuted] at (0,-2.05) {#4};
  \node[utplabel,anchor=north,align=center] at (0,-2.42) {#5};
 \end{scope}}

\ttree{0}{Inorder}{n4/1,n2/2,n5/3,n1/4,n3/5}{left, \textbf{root}, right}{D B E A C}
\ttree{5.6}{Preorder}{n1/1,n2/2,n4/3,n5/4,n3/5}{\textbf{root}, left, right}{A B D E C}
\ttree{11.2}{Postorder}{n4/1,n5/2,n2/3,n3/4,n1/5}{left, right, \textbf{root}}{D E B C A}
\end{tikzpicture}
\end{labgraphic}
\end{utpcompactlabfigure}


\lstinputlisting[firstline=31,lastline=62]{all_experiments/Lab07/binary_tree.cpp}


\tcbbreak

\begin{utpsteps}\setcounter{enumi}{2}\item \textbf{Implement deleteValue, destroy.}\end{utpsteps}

\texttt{binary\_tree.cpp} --- lines 63--97.

% Worked lab input; chapter fonts, colours and TikZ styles retained.
\begin{utpcompactlabfigure}[General binary-tree deletion: copy the last level-order value, then remove its leaf]
\begin{labgraphic}
\begin{tikzpicture}[font=\small]
\begin{scope}[xshift=0cm,yshift=0cm]
\node[utplabel,anchor=south,color=UTPNavy] at (0,.65) {Updated 6 to 7; inserted 10};
\node[utpnode,minimum size=7mm] (n0) at (0,0) {8};
\node[utpnode,minimum size=7mm] (n1) at (-1.4,-1) {4};
\node[utpnode,minimum size=7mm] (n2) at (1.4,-1) {12};
\node[utpnode,minimum size=7mm] (n3) at (-2.1,-2) {2};
\node[utpnodehot,minimum size=7mm] (n4) at (-0.7,-2) {7};
\node[utpnodehot,minimum size=7mm] (n5) at (0.7,-2) {10};
\draw[utpedge] (n0)--(n1);
\draw[utpedge] (n0)--(n2);
\draw[utpedge] (n1)--(n3);
\draw[utpedge] (n1)--(n4);
\draw[utpedge] (n2)--(n5);
\end{scope}\begin{scope}[xshift=7cm,yshift=0cm]
\node[utplabel,anchor=south,color=UTPNavy] at (0,.65) {Deleted the value 4};
\node[utpnode,minimum size=7mm] (n0) at (0,0) {8};
\node[utpnodehot,minimum size=7mm] (n1) at (-1.4,-1) {10};
\node[utpnode,minimum size=7mm] (n2) at (1.4,-1) {12};
\node[utpnode,minimum size=7mm] (n3) at (-2.1,-2) {2};
\node[utpnode,minimum size=7mm] (n4) at (-0.7,-2) {7};
\draw[utpedge] (n0)--(n1);
\draw[utpedge] (n0)--(n2);
\draw[utpedge] (n1)--(n3);
\draw[utpedge] (n1)--(n4);
\end{scope}
\draw[utparrow] (2.5,-.8)--(4.3,-.8);
\node[utplabel,align=center] at (3.5,-3) {\texttt{target->data = deepest->data}; clear the parent's link; free deepest\\Array: $[8,4,12,2,7,10]\longrightarrow[8,10,12,2,7,\text{empty}]$};
\end{tikzpicture}
\end{labgraphic}
\end{utpcompactlabfigure}


\lstinputlisting[firstline=63,lastline=97]{all_experiments/Lab07/binary_tree.cpp}


\tcbbreak

\begin{utpsteps}\setcounter{enumi}{3}\item \textbf{Implement present, arrayInorder, arrayPreorder and helpers.}\end{utpsteps}

\texttt{binary\_tree.cpp} --- lines 98--124.

% Adapted from the chapter tree-array drawing: five-node lab starting state.
% Keep the chapter drawing, type sizes, colours, and node styles unchanged.
\begin{adslabfigure}[A binary tree stored in an array]
\begin{adslabgraphic}
\begin{tikzpicture}[font=\small]
% tree
\node[utpnode,minimum size=8mm] (v0) at (2.55,1.5) {8};
\node[utpnode,minimum size=8mm] (v1) at (1.25,0.45) {4};
\node[utpnode,minimum size=8mm] (v2) at (3.85,0.45) {12};
\node[utpnode,minimum size=8mm] (v3) at (0.6,-0.6) {2};
\node[utpnode,minimum size=8mm] (v4) at (1.9,-0.6) {6};
\draw[utpedge] (v0)--(v1); \draw[utpedge] (v0)--(v2);
\draw[utpedge] (v1)--(v3); \draw[utpedge] (v1)--(v4);

\foreach \n/\idx in {v0/0,v1/1,v2/2,v3/3,v4/4}{
  \node[circle,fill=UTPBlue,text=white,font=\tiny\bfseries,inner sep=1.1pt]
    at ($(\n.north east)+(0.06,0.06)$) {\idx};
}
% array
\begin{scope}[yshift=-2.35cm]
\foreach \v [count=\i from 0] in {8,4,12,2,6}{
  \node[utpcell,minimum width=1.0cm,minimum height=0.82cm] (g\i) at (1.05*\i,0) {\v};
  \node[circle,fill=UTPBlue,text=white,font=\tiny\bfseries,inner sep=1.1pt] (arrayindex\i)
    at (1.05*\i,0.55) {\i};
}
\node[utplabel,anchor=east] at (-0.65,0) {array};
\end{scope}
% Route around the children, then enter the mapped cell through the gap.
\draw[utpptr,draw=UTPOrange,dashed]
  (v1.east)
  .. controls (2.15,0.45) and (2.55,0.22) .. (2.55,-0.35)
  .. controls (2.55,-0.95) and (2.55,-1.38) .. (2.55,-1.5)
  .. controls (2.35,-1.62) and (1.9,-1.55) .. (1.575,-1.55)
  -- (arrayindex1.center);
\node[utptag,anchor=west,align=left] at (6.6,-0.1)
  {index 1 holds 4\\children: $2(1)+1=3$\\\phantom{children: }$2(1)+2=4$\\$\rightarrow$ 2 and 6\ \checkmark};
\end{tikzpicture}
\end{adslabgraphic}
\end{adslabfigure}


\lstinputlisting[firstline=98,lastline=124]{all_experiments/Lab07/binary_tree.cpp}


\tcbbreak

\begin{utpsteps}\setcounter{enumi}{4}\item \textbf{Implement arrayDeleteValue, printIndexTable, sameShape.}\end{utpsteps}

\texttt{binary\_tree.cpp} --- lines 125--162.

\lstinputlisting[firstline=125,lastline=162]{all_experiments/Lab07/binary_tree.cpp}


\tcbbreak

\begin{utpsteps}\setcounter{enumi}{5}\item \textbf{Connect the test cases in main.}\end{utpsteps}

\texttt{binary\_tree.cpp} --- lines 163--199.

\textbf{Trace checkpoint for the supplied main.} After inserting 10 as the left child of 12, its array position is 2 times 2 plus 1 = 5. Deleting value 4 uses the deepest rightmost node, 10, as replacement in this general binary-tree exercise. That is different from the BST successor rule in the next lab. Check linked and array traversals after every mutation, not only their initial shape.

\lstinputlisting[firstline=163,lastline=199]{all_experiments/Lab07/binary_tree.cpp}


\tcbbreak

\begin{utpsteps}\setcounter{enumi}{6}\item \textbf{Connect the test cases in main.}\end{utpsteps}

\texttt{binary\_tree.cpp} --- lines 200--240. Continue in the same file; do not add another header or opening brace.

\lstinputlisting[firstline=200,lastline=240]{all_experiments/Lab07/binary_tree.cpp}


\tcbbreak

\textbf{Build, run, and inspect binary\_tree.}

From this program\textquotesingle s folder, run these commands in PowerShell:

\begin{lstlisting}[style=utpterminal,language={}]

g++ -std=c++17 -Wall -Wextra binary_tree.cpp -o binary_tree.exe

.\binary_tree.exe

\end{lstlisting}

\textbf{Expected output excerpt.} Compare values, order, and pass/fail messages; spacing need not match.

\begin{lstlisting}[style=utpterminal,language={}]

== Predicted before running (write these down first) ==
  inorder   2 4 6 8 12
  preorder  8 4 2 6 12
  postorder 2 6 4 12 8

== Traversals ==
  inorder   linked: 2 4 6 8 12   array: 2 4 6 8 12 
  preorder  linked: 8 4 2 6 12   array: 8 4 2 6 12 
  postorder linked: 2 6 4 12 8   array: 2 6 4 12 8 
  [ok]   both representations describe the same shape
  [ok]   all three traversals agree between the representations
  [ok]   the predictions match
  [ok]   each traversal visits every stored node exactly once

== Array index table ==

\end{lstlisting}

The complete captured run is in \texttt{labs/outputs/binary\_tree.txt}.

\begin{tryit}

Delete a missing value and then the only node of a fresh tree in both representations. Compare traversal outputs and shape, and confirm all linked allocations are released.

\end{tryit}
